\subsubsection{与 K 理论的关系}\label{dolfo}

我们以一种相当技术性的方式求得了线丛 $\L\to D$ 的次数 $c_0$，但实际上，这个答案（式 (\ref{ywo})）有一个很好的解释，文献 \cite{FW} 第 4.3 节实质上已经描述了这一解释.

让我们把 $W=T^*C$ 看成实辛流形，其辛形式取为全纯辛形式 $\Omega_I$ 的虚部（或实部）. 我们考虑具有这一辛形式的 $W$ 的 $A$-模型. 我们将 $C\subset W$ 看成拉格朗日子流形，并考虑支撑在 $C$ 上的 $A$-膜.

直观地说，支撑在 $C$ 上的秩为 1 的 $A$-膜带有一个平坦的陈-帕顿线丛\index{chan-paton-line-bundle@陈-帕顿线丛} $\U\to C$. 但实际上，$D$-膜的 $K$ 理论解释意味着存在一种扭曲，它涉及 $C$ 在 $W$ 中的法丛的平方根. 其结果是，$\U$ 应当是 $C$ 上的一个平坦 ${\mathrm{Spin}_c}$ 结构，而非平坦线丛. 用复几何的语言说，这意味着 $\U$ 的次数不应为零，而应等于典范丛 $K\to C$ 的一个平方根的次数，也就是说，它应当是 $g_C-1$.\transnote{这里“平坦 ${\mathrm{Spin}_c}$ 结构”的来历是 D-膜的 $K$ 理论分类：包裹子流形 $C$ 的 D-膜，其陈-帕顿丛必须与 $C$ 在 $W$ 中法丛的旋量丛张量在一起才有意义，而 $W=T^*C$ 的法丛同构于余切丛 $K$，故要求 $\U^2\cong K$。当 $w_2(C)\neq0$ 时 $K$ 没有整体平方根，只能存在平坦 ${\mathrm{Spin}_c}$ 结构，其次数为 $g_C-1$。（来源：A 级，Witten，\emph{D-Branes And K-Theory}，arXiv:hep-th/9810188 §4.3“The Spin\_c Case”，判据 $c_1(L)\equiv w_2(N)\bmod 2$；原文 §2.6.4 与式 (\ref{tywo})，2026-11-19 核对.）}

类似地，支撑在 $C$ 上的秩为 $N$ 的 $A$-膜的陈-帕顿丛\index{chan-paton-bundle@陈-帕顿丛} $E$ 的次数应为 $N(g_C-1)$. 这意味着，如果我们希望线丛 $\L\to D$ 具有如下性质：$E=\psi_*(\L)$ 是支撑在 $C$ 上的一个 $A$-膜的陈-帕顿丛，那么它的次数就必须不是式 (\ref{ywo}) 中定义的 $c_0$，而是
\begin{equation}\label{tywo}c_0'=c_0+N(g_C-1)=N^2(g_C-1)=g_D-1. \end{equation}
但 $D$ 也是 $W$ 的一个拉格朗日子流形，而 $\L\to D$ 的次数取这个值，恰好意味着赋予了陈-帕顿丛 $\L$ 的 $D$ 可以解释为一个 $A$-膜.

因此，我们可以把希格斯丛和谱覆盖的构造看成在支撑于 $D$ 上的秩为 1 的 $A$-膜与支撑于 $C$ 上的秩为 $N$ 的 $A$-膜之间给出了一种对应.

\subsubsection{酉群}
\label{unitary}

我们可以用类似的方式，分析相关群 $U(N)$ 和 $PSU(N)=SU(N)/\Bbb{Z}_N$ 的希钦纤维化的纤维.

首先考虑 $G=U(N)$. 对于这个例子，我们应当修改开头的构造，不再要求 $\varphi$ 无迹. 相应地，在定义 ${\cal V}$ 时，我们在高次迹之外还包括 $\Tr\,\varphi$，结果 ${\cal V}$ 的维数为 $(g_C-1)N^2+1$. 此外，对于 $U(N)$，我们不会把 $E$ 的行列式限制为平凡的，因此 $(g_C-1)N^2+1$ 正是纤维 $\FF$ 应有的维数. 这些维数彼此相等，并且等于 $\frac{1}{2}\dim\,\MH$，后者对于 $G=U(N)$ 为 $N^2(g_C-1)+1$. 由于 $E$ 的行列式不受平凡性约束，$\psi_*({\cal L})$ 可以是任意的，粗略地说，${\cal L}$ 可以是 $D$ 上的任意线丛. 因而，$U(N)$ 的希钦纤维化的纤维就是参数化 $D$ 上线丛的皮卡群 ${\rm Pic}(D)$.

更精确地说，我们需要注意到，这个问题有两个相关的版本. 对于 $U(N)$，希格斯丛 $(E,\varphi)$ 有一个整数值拓扑不变量，即 $E$ 的第一陈类，记为 $c_1(E)$. 稳定希格斯丛的第一陈类可以取任意值. 对于这个不变量的任意取值，希钦纤维化都有意义. 当 $c_1(E)=m$ 时，希钦纤维化的纤维是模空间 ${\rm Pic}_{c_0+m}(D)$，它参数化 $D$ 上次数为 $c_0+m$ 的线丛，其中 $c_0$ 由 (\ref{ywo}) 给出. 这可由与之前相同的分析得出.

然而，稳定希格斯丛与希钦方程解之间的等价仅在第一陈类为零时成立. 希钦方程之一 $F-\phi\wedge \phi=0$ 取迹后化为 $\Tr\,F=0$，这意味着第一陈类必须为零. 因此，希钦方程解的模空间 $\MH(U(N),C)$（或我们场论的经典真空）实际上对应于稳定希格斯丛模空间中 $c=0$ 的连通分支. 所以，$\MH(U(N),C)$ 的希钦纤维化的纤维为
\begin{equation}\label{picun}\FF_{U(N)}={\rm
Pic}_{c_0}(D).\end{equation} 与我们上面对 $SU(N)$ 的讨论完全相同，当 $N$ 为奇数时，这个纤维典范地是一个阿贝尔簇，但当 $N$ 为偶数时，它还不能完全算作阿贝尔簇. 确切地说，一旦在 $C$ 上选定自旋结构\index{spin-structure@自旋结构} $K^{1/2}$，它就成为阿贝尔簇.\transnote{“典范地成为阿贝尔簇”指这个复环面有一个不依赖额外选择的零点，从而可等同于 $\mathrm{Pic}^0$ 型的对象。$N$ 为奇数时，典范点 $\psi^*(K^{(N-1)/2})$ 由 $K$ 的整数次幂直接给出；$N$ 为偶数时 $(N-1)/2$ 不是整数，必须先取定 $K^{1/2}$（自旋结构），典范点才存在，这正是“还不能完全算作阿贝尔簇”的意思。后文 (2.38) 中除以 $\Gamma_N$ 恰好消去对自旋结构的依赖。（来源：A 级，原文 §2.6.3 与 §2.6.5 式 (2.36)–(2.38)；A 级，Frenkel–Witten，arXiv:0710.5939 §10 对相关修正的讨论；2026-11-19 核对.）}

对于 $U(N)$，我们主要关注次数为零的希格斯丛 $(E,\varphi)$，因为它们表示我们的 $\sigma$-模型的经典真空模空间 $\MH(U(N),C)$ 中的点. 尽管如此，我们仍可问：其他次数的希格斯丛 $(E,\varphi)$ 是否与几何朗兰兹有关. 答案是，如果所问的问题不能充分地通过把四维规范理论约化为靶空间是 $\MH$ 的二维 $\sigma$-模型来处理，那么它们就可能有关. 更多讨论见第 \ref{elmag} 节.


\subsubsection{群 $PSU(N)$}
\label{psun}

最后，我们来确定群 $PSU(N)=SU(N)/\Bbb{Z}_N$ 的希钦纤维化的纤维，其中 $\Bbb{Z}_N$ 是 $SU(N)$ 的中心.

$C$ 上的 $PSU(N)$ 丛可以在拓扑上非平凡. 先考虑拓扑平凡的情形. $C$ 上结构群为 $PSU(N)$ 的拓扑平凡丛 $E'$ 可以提升为 $C$ 上的 $SU(N)$ 丛 $E$，但这种提升并不完全唯一. 唯一性之所以失效，是因为可以用一个阶为 $N$ 的线丛 ${\cal S}$ 扭曲 $E$，或者等价地，用一个单值变换取值于 $SU(N)$ 中心的平坦线丛扭曲它. 这样一个阶为 $N$ 的平坦丛在投影到 $PSU(N)$ 时就消失了. 设 $\Gamma_N$ 是 $C$ 上阶为 $N$ 的线丛所成的群. 那么，对于拓扑平凡的 $PSU(N)$ 丛，希钦纤维化的纤维为
\begin{equation}
\label{mormot}
\FF^{(0)}_{PSU(N)}=\FF_{SU(N)}/\Gamma_N=\Lambda^{-1}({\cal
O})/\Gamma_N.
\end{equation}
（元素 ${\cal S}\in \Gamma_N$ 通过 ${\cal L}\to {\cal L}\otimes
\psi^*({\cal S})$ 作用于 ${\cal L}\in
\Lambda^{-1}({\cal O})$.）$\FF^{(0)}_{PSU(N)}$ 始终典范地是一个阿贝尔簇. 对于奇数 $N$，即使在除以 $\Gamma_N$ 之前，$\FF_{SU(N)}$ 就已经是阿贝尔簇. 对于偶数 $N$，为了使 $\FF_{SU(N)}$ 成为阿贝尔簇，我们必须在 $C$ 上选定一个自旋结构. 然而，除以 $\Gamma_N$ 恰好消除了对自旋结构的依赖. 事实上，改变自旋结构相当于用一个阶为 2 的线丛来扭曲，而对于偶数 $N$，这样的线丛是 $\Gamma_N$ 的元素.

注意到 $PSU(N)$ 与 $PU(N)=U(N)/U(1)$ 相同，我们就能找到描述 $\FF^{(0)}_{PSU(N)}$ 的另一种有用方式. 对于 $C$ 上次数为零的线丛 ${\cal S}$，$\psi^*({\cal S})$ 是 $D$ 上次数为零的线丛. 因而，对于 ${\cal L}\in
\FF_{U(N)}= {\rm Pic}_{c_0}(D)$，张量积 ${\cal L}\otimes
\psi^*({\cal S})$ 也属于 ${\rm Pic}_{c_0}(D)$. 这给出 ${\rm Jac}(C)$ 在 ${\rm Pic}_{c_0}(D)$ 上的群作用，我们将其商记为 ${\rm Pic}_{c_0}(D)/{\rm Jac}(C)$. 那么
\begin{equation}\label{publy} \FF^{(0)}_{PSU(N)}={\rm
Pic}_{c_0}(D)/{\rm Jac}(C). \end{equation} 这一断言的依据是，由于 $\psi_*({\cal L}\otimes \psi^*{\cal
S}))=\psi_*({\cal L})\otimes {\cal S}$（根据 (\ref{xofo})），并且 $\psi_*({\cal L})$ 的秩为 $N$，我们有 $\det(\psi_*({\cal L}\otimes \psi^*({\cal S})))=\psi_*({\cal
L})\otimes {\cal S}^N$. 因此，可以选取 ${\cal S}$ 使得 $\det(\psi_*({\cal L}\otimes \psi^*({\cal S})))={\cal O}$，从而 ${\cal L}\otimes \psi^*({\cal S})$ 成为 $\FF_{SU(N)}$ 的元素. 这个条件并未完全唯一地确定 ${\cal S}$. 我们仍然可以用 $\Gamma_N$ 的一个元素来扭曲 ${\cal S}$. 考虑到这一点，就可以看出 (\ref{mormot}) 与 (\ref{publy}) 等价.

对于拓扑平凡情形的这第二种描述，是分析拓扑非平凡 $PSU(N)$ 丛的希钦纤维化的一个有用出发点. 从拓扑上说，$PSU(N)$ 丛 $E$ 由一个特征类\index{characteristic-class@特征类} $\xi$ 分类，其取值于 $H^2(C,\Bbb{Z}_N)\cong \Bbb{Z}_N$.（对于 $N=2$，$\xi$ 是第二 Stiefel–Whitney 类 $w_2(E)$.）这个特征类度量将 $PSU(N)$ 丛“提升”到 $SU(N)$ 的障碍. 设 $\MH^{(d)}$ 是 $\MH(PSU(N),C)$ 中参数化特征类为 $\xi=d$ 的希格斯丛的连通分支.\transnote{这类特征类就是所谓 \emph{'t Hooft 磁通}：$C$ 是二维的，$PSU(N)=SU(N)/\Bbb Z_N$ 丛的分类只依赖 $\pi_1(PSU(N))=\Bbb Z_N$，故 $H^2(C,\Bbb Z_N)\cong\Bbb Z_N$；$N=2$ 时 $PSU(2)=SO(3)$，$\xi$ 就是 $w_2(E)$。若 $PSU(N)$ 丛来自 $U(N)$ 丛 $E$，则 $\xi\equiv c_1(E)\bmod N$，这解释了下一段的提升判据。（来源：A 级，[KW]（arXiv:hep-th/0604151）§7.1 关于 $G_{ad}$ 丛特征类 $\xi(E)\in H^2(M,\Bbb Z)$ 与 $SO(3)$ 情形即 $w_2$ 的讨论；原文 §2.6.6；2026-11-19 核对.）}

尽管拓扑非平凡的 $PSU(N)$ 丛不能提升到 $SU(N)$，它却可以非唯一地提升到 $U(N)$. 事实上，如果 $E$ 是任意 $U(N)$ 丛，那么相应的伴随丛 ${\rm ad}(E)$ 是 $PSU(N)$ 丛，而且每个 $PSU(N)$ 丛都可以由某个 $E$ 以这种方式得到. 如果 $E$ 的第一陈类是 $d$，或者更一般地，它模 $N$ 同余于 $d$，那么 $PSU(N)$ 丛 ${\rm ad}(E)$ 的特征类为 $\xi=d$.

如果 $E_1$ 和 $E_2$ 是两个不同的 $U(N)$ 丛，那么相应的 $PSU(N)$ 丛 ${\rm ad}(E_1)$ 和 ${\rm ad}(E_2)$ 同构，当且仅当对某个线丛 ${\cal S}\to C$ 有 $E_1\cong E_2\otimes{\cal S}$.（注意，在这一操作中，$E_1$ 和 $E_2$ 的第一陈类模 $N$ 相等，因为 $c_1(E_1)=c_1(E_2)+Nc_1({\cal S})$. 它们模 $N$ 约化为 $d=\xi({\rm ad}(E))$.）这导出 (\ref{publy}) 的一个类似式. 对于 $\xi=d$ 的 $PSU(N)$ 丛，希钦纤维 $\FF^{(d)}_{PSU(N)}$ 为
\begin{equation}\label{blonko} \FF^{(d)}_{PSU(N)}={\rm
Pic}_{c_0+d}(D)/{\rm Jac}(C).\end{equation} 这里的思想是，${\rm Pic}_{c+d}(D)$ 中的一个点确定 $C$ 上一个满足 $c_1(E)=d$ 的 $U(N)$ 丛 $E$；${\cal S}\in {\rm Jac}(C)$ 的作用产生 $E\to E\otimes {\cal S}$ 的效果，而不影响相应的 $PSU(N)$ 丛 ${\rm ad}(E)$.

还有一个有用的、虽然看起来较不典范的 (\ref{mormot}) 的类似式. 给定一个我们记为 ${\rm ad}(E)$ 的 $PSU(N)$ 丛，由于存在 $E\to E\otimes {\cal S}$ 的自由度，它到 $U(N)$ 的提升 $E$ 远非唯一. 该作用在 $\det(E)$ 上为 $\det(E)\to \det(E)\otimes {\cal S}^N$. 因此，只考虑次数为零的 ${\cal S}$，我们可以在 $\det(E)$ 所属的 ${\rm Pic}(C)$ 的连通分支内任意调整它，随后仍然可以用一个阶为 $N$ 的线丛，即 $\Gamma_N$ 的一个元素，来扭曲 $E$. 我们可以在 $C$ 上选取任意方便的、次数为 $d=\xi({\rm ad}(E))$ 的线丛 ${\cal S}_0$，并要求 $\det(E)={\cal S}_0$. 例如，可以选取一点 $\pp\in C$，并取 ${\cal S}_0={\cal O}(\pp)^d$. 这样就得到 (\ref{mormot}) 的类似式：
\begin{equation}
\label{normat} \FF^{(d)}_{PSU(N)}=\Lambda^{-1}({\cal S}_0)/\Gamma_N,
\end{equation}
它表达了如下事实：通过 $J(C)$ 的作用，我们可以固定 $E$ 的行列式，而之后仍可施加 $\Gamma_N$ 的作用.

最后一点说明是，我们在 (\ref{mormot}) 和 (\ref{normat}) 中确定的，恰好是 $\MH(PSU(N),C)$ 的希钦纤维化的纤维；该模空间定义为所有希钦方程的解的空间除以所有取值于 $PSU(N)$ 的规范变换所成的群. 有时，在 $\MH(PSU(N),C)$ 的万有覆盖上工作也很方便；只除以能连续形变到恒等变换的规范变换子群，就能得到这一覆盖. 在希钦纤维化的纤维上，这恰好相当于在上述构造中不除以 $\Gamma_N$. 因此，如果 $\hat
\MH^{(d)}(PSU(N),C)$ 是仅除以连通于恒等变换的规范变换所得到的 $\MH^{(d)}(PSU(N),C)$ 的覆盖，那么相应的希钦纤维化的纤维为
\begin{equation}\label{palooko} \hat \FF^{(d)}_{PSU(N)}=\Lambda^{-1}({\cal
L}_0).\end{equation} 如果取 $d=0$，$\hat\MH^{(d)}(PSU(N),C)$ 就化为 $\MH(SU(N),C)$. 事实上，对于 $d=0$，可以取 ${\cal L}_0={\cal O}$，而 $\hat \FF^{(d)}_{PSU(N)}$ 化为 $\FF_{SU(N)}$.

人们可能会问，$\MH(PSU(N))$ 的这些拓扑非平凡连通分支的对偶是什么. 这里相关的对偶性是 $\MH(PSU(N))$ 与 $\MH(SU(N))$ 之间的 $S$-对偶性或镜像对称\index{mirror-symmetry@镜像对称}，它是这些规范群的几何朗兰兹对偶性的基础；完整解释见 \cite{KW}. 文献 \cite{HauselThaddeus} 从数学角度解释了，$\MH^{(d)}(PSU(N))$ 的对偶就是我们熟悉的 $\MH(SU(N))$，但要赋予它某个非平凡的 gerbe. 下文第 \ref{gerbes} 节将详细阐述这里相关的 gerbe 概念. 另见 \cite{KW} 第 7 节，其中把有关的 gerbe 解释为 $\sigma$-模型的离散 $B$-场，并将其与电磁对偶性\index{electromagnetic-duality@电磁对偶性}的一个离散版本联系起来. 关于任意 $G$ 下这些问题的分析，见 \cite{donga}.\transnote{这里“非平凡的 gerbe”可理解为 $\sigma$-模型的离散 $B$-场：$U(1)$-gerbe 的同构类由 $H^2(X,U(1))$ 分类，其上的向量丛只能“扭曲地”存在，即转移函数只须满足 $v_{ki}v_{jk}v_{ij}=b_{ijk}$。$PSU(N)$ 分支与 $SU(N)$ 分支对偶时，两边相差的正是这份离散数据。（来源：A 级，[KW]（arXiv:hep-th/0604151）§7.2 关于平坦 $B$-场 $b\in H^2(X,U(1))$ 与 twisted vector bundle 的讨论；B 级，nLab“gerbe”\url{https://ncatlab.org/nlab/show/gerbe}（Last revised 2026-02-14，2026-11-19 访问）给出 gerbe 的叠定义，待核；原文 §2.6.6 与 §3.2.）}

\subsubsection{其他规范群的谱覆盖}\label{spectrother}

我们仅对 $G=SU(N)$ 以及与之密切相关的群 $U(N)$ 和 $PSU(N)$ 描述了谱覆盖理论. 对于其他 $G$ 会怎样？

如果 $G$ 有一个方便的“小”表示，就可以发展出与上述描述有些相似的理论. 对于 $SO(N)$、$Sp(N)$，甚至 $G_2$，这都是切实可行的 \cite{HitchinDuke, KP, Hthree}.

一般而言——例如对于 $G=E_8$——不存在方便的“小”表示，需要使用更抽象的方法. 见 \cite{donga}.\transnote{对一般 $G$，“谱覆盖”要换成 \emph{cameral 覆盖}：它是 $C$ 的一个 $|W|$ 重覆盖（$W$ 为 Weyl 群），在正则点处的纤维是 $G$ 的 Cartan 子代数。谱覆盖只对 $GL(N)$ 或具有“小”表示的群才有简单意义，cameral 覆盖则对所有 $G$ 都有定义；两者的关系是：谱覆盖为 cameral 覆盖与小表示所做的纤维积。（来源：A 级，Donagi–Gaitsgory，\emph{The Gerbe Of Higgs Bundles}，arXiv:math/0005132 §2“Higgs bundles and cameral covers”与 §9“Spectral covers versus cameral covers for $G=GL(n)$”；原文 §2.6.7；2026-11-19 核对.）}

\subsection{特选截面}
\label{distinguished}
\subsubsection{$SU(N)$ 的情形}

为了熟悉这些思想，同时也因为这确实有用，我们将更仔细地考察某一族特选希钦对 $(E,\varphi)$.

从 $G=SU(2)$ 开始，取 $E$ 为秩二复向量丛 $E=K^{1/2}\oplus K^{-1/2}$. 这个丛不稳定；它包含次数为正的线子丛 $K^{1/2}$. 然而，仍然存在稳定的希钦对 $(E,\varphi)$，我们将遵循 \cite{Hitchin} 对它们进行描述.

$E$ 不稳定这一事实的一个表现是，与稳定丛不同，它具有非常值的自同构. 如果把 $E$ 的一个截面写为 $\begin{pmatrix}s\\ t\end{pmatrix}$，其中 $s$ 是 $K^{1/2}$ 的截面，$t$ 是 $K^{-1/2}$ 的截面，那么 $E$ 的 $SL(2,\C)$ 自同构作用为
\begin{equation}
\label{uti}
\begin{pmatrix}\lambda & \tau\\
0&\lambda^{-1}\end{pmatrix}.
\end{equation}
这里 $\lambda$ 是非零复数，$\tau$ 是 $H^0(C,K)$ 的元素. 我们希望在这些自同构对 $\varphi$ 的作用下，对稳定希钦对 $(E,\varphi)$ 进行分类.

$H^0(C,K\otimes\ad(E))$ 的一般全纯截面 $\varphi$ 可以写为
\begin{equation}
\label{buti}
\begin{pmatrix} v &\; w/2 \\ u &\;
-v\end{pmatrix},
\end{equation}
其中 $u$ 是复数，$v$ 是 $H^0(C,K)$ 的元素，$w$ 是 $H^0(C,K^2)$ 的元素.（乘在 $w$ 上的因子 1/2 是为了方便.）要使对 $(E,\varphi)$ 稳定，$u$ 必须非零，因为若 $u=0$，则次数为正的线丛 $K^{1/2}\subset E$ 在 $\varphi$ 下不变. 如果 $u\not=0$，我们可以选取 $\lambda$ 和 $\tau$，使得 $u=1$、$v=0$（这种选择模去 $SL(2,\C)$ 的中心 $\{\pm 1\}$ 后是唯一的），此时
\begin{equation}
\label{hgy}
\varphi = \begin{pmatrix} 0 & w/2 \\ 1 & 0\end{pmatrix}.
\end{equation}
如 \cite{Hitchin} 所证明的，任何这样的对 $(E,\varphi)$ 都是稳定的（而且 $E$ 是仍存在这种稳定对的最不稳定的丛）. 显然 $\Tr\,\varphi^2=w$，因此，对于这个特定的 $E$，给定 $\Tr\,\varphi^2$ 后，模去自同构存在唯一的稳定对 $(E,\varphi)$. 所以，具有这个 $E$ 的稳定希钦对所成的族 $V_E$ 与希钦纤维化的每个纤维 $\FF$ 恰好交于一点. 因而，在选定 $K^{1/2}$ 后，族 $V_E$ 给出的截面
\begin{equation}\label{zelo}\zeta:\V\to\MH\end{equation}
是希钦纤维化的一个特选截面\index{distinguished-section@特选截面}. 正如我们在第 \ref{hitchfibr} 节末尾所解释的，对于一般的稳定 $E$，$V_E$ 与 $\FF$ 的交点数为 $2^{3g-3}$.\transnote{“特选截面”指每个希钦纤维里都有一个由不稳定丛给出的标准点：$SL(2,\Bbb C)$ 情形取 $E=K^{1/2}\oplus K^{-1/2}$，其上的希格斯场在 $w\neq0$ 时稳定，于是纤维化得到一个不依赖参数的标准截面；而对一般稳定 $E$，截面与纤维的交点数为 $2^{3g-3}$，说明这种“唯一标准点”只在退化（不稳定）丛处成立。（来源：A 级，Hitchin，\emph{The self-duality equations on a Riemann surface}，Proc. London Math. Soc. (3) 55 (1987) 59–126；原文 §2.7.1 与 §2.6.3；2026-11-19 核对.）}

谱曲线 $D$ 由方程 $\det(y-\varphi)=0$ 或 $y^2=w/2$ 定义. 通过求解本征向量方程 $\varphi\uppsi=y\uppsi$，即
\begin{equation}\label{ivox}
\begin{pmatrix} 0 & w/2 \\ 1 & 0\end{pmatrix}\uppsi = y
\uppsi.
\end{equation}
可以求得线丛 ${\cal N}\to D$. 局部地，可以方便地取
\begin{equation}
\label{yupo}
\uppsi=\begin{pmatrix} y\alpha \\ \alpha\end{pmatrix},
\end{equation}
来求解，其中 $\alpha$ 是 $K^{-1/2}$ 的任意局部全纯截面. 这给出了从 $K^{-1/2}$ 的截面到 ${\cal N}$ 的截面的自然映射，从而在这个例子中，${\cal N}$ 与 $\psi^*(K^{-1/2})$ 相同.

$\varphi-y$ 的“余核”定义了另一个线丛 ${\cal L}$，根据 (\ref{roygo})，它为 ${\cal L}={\cal N}\otimes
\psi^*(K)=\psi^*(K^{1/2})$. 丛 $E$ 应当是 $\psi_*({\cal L})$，根据 (\ref{xofo})，后者为 $K^{1/2}\oplus K^{-1/2}$. 这确实同构于我们开始时使用的丛 $E$.

对于 $G=SU(N)$，与此类似的是取 $E=K^{(N-1)/2}\oplus K^{(N-1)/2-1}\oplus\dots\oplus
K^{-(N-1)/2}$. 我们将 $E$ 的一个截面表示为列向量
\begin{equation}
\label{plook}
\begin{pmatrix}s_1\\
s_2\\ \vdots\\ s_N\end{pmatrix},
\end{equation}
其中 $s_i$ 是 $K^{(N+1)/2-i}$ 的截面. 与上述类似的分析表明，具有这个 $E$ 的稳定希钦对 $(E,\varphi)$ 可以写成如下形式
\begin{equation}
\label{plyo}
\varphi=\begin{pmatrix}0 & w_2/2 & w_3/3 & \dots
&w_{N-1}/(N-1)& w_N/N\\ 1 & 0 & 0 & \dots &0& 0 \\ 0 & 1 & 0 &
\dots & 0&0\\ \vdots & \vdots &\vdots & \vdots& \ddots &\vdots\\
0&0&0&\dots &1&0\end{pmatrix},
\end{equation}
其中紧靠主对角线下方的元素为 1，$w_n$、$n=2,\dots,N$ 取值于 $H^0(C,K^n)$. 我们有 $\Tr\,\varphi^n=w_n$. 因而，对于这个 $E$，在希钦纤维化的每个纤维 $\FF$ 上再次存在唯一的稳定希钦对，从而给出希钦纤维化的一个自然截面. 与在 $SU(2)$ 情形下的发现类似，通过求解本征向量方程，可以得到 ${\cal
N}\cong \psi^*(K^{-(N-1)/2})$，进而利用 (\ref{roygo}) 得到 ${\cal L}\cong \psi^*(K^{(N-1)/2})$. 我们预期可由 $\psi_*({\cal L})$ 恢复 $E$. 事实上，从 (\ref{xofo}) 确实得到 $\psi_*({\cal L})=K^{(N-1)/2}\oplus
K^{(N-1)/2-1}\oplus \dots\oplus K^{-(N-1)/2}=E$.

\subsubsection{任意 $G$ 的希钦纤维化的截面}

所有这些对于任意单李群 $G$ 都有类似的版本，不过我们并不打算给出完整解释.


一般情形的出发点是利用 Kostant 的主 $SU(2)$ 嵌入，将 $SU(2)$ 丛 $K^{1/2}\oplus K^{-1/2}$ 嵌入 $G$. 这给出一个特殊的不稳定 $G$-丛 $E$. 对于这个 $E$，在希钦纤维化的任意给定纤维上，存在唯一的稳定希钦对 $(E,\varphi)$. 文献 \cite{BeD} 利用“oper”理论，以略有不同的表述证明了这一点. 因而，这一构造对于任意 $G$ 都给出希钦纤维化的一个几乎典范的截面. 我们称这个截面几乎典范，是因为它一般可能依赖于 $K^{1/2}$ 的选择.\transnote{“oper”是 Beilinson–Drinfeld 引入的概念：粗略地说，它是带有一个满足额外横截与正则性条件的联络的 $G$-丛（$\mathrm{PSL}(2)$ 情形下，oper 等同于 $C$ 上的射影结构）。oper 的模空间给出希钦纤维化的一个截面，其取值正是 Kostant 主 $SU(2)$ 嵌入所确定的不稳定丛，故截面不依赖纤维化参数，只因 $K^{1/2}$ 的选择而“几乎”典范。（来源：A 级，Beilinson–Drinfeld，\emph{Quantization Of Hitchin's Integrable System And Hecke Eigensheaves}，preprint，即原文 §2.7.2 所引 [12]，opers 概念与 Hitchin 系统的关系见该书导言；原文 §2.7.2；2026-11-19 核对.）}

事实上，对于任意 $G$，Kostant 嵌入都是将 $SU(2)$ 的李代数嵌入 $G$ 的李代数的嵌入. 在群的层面上，嵌入 $G$ 的可以是 $SU(2)$，也可以是 $PSU(2)=SO(3)$. 在后一种情形下，作 Kostant 嵌入时，对 $K^{1/2}$ 选择的依赖就消失了，我们得到希钦纤维化的一个真正典范的截面. 在前一种情形下，对于 $K^{1/2}$ 的每一种选择，都得到希钦纤维化的一个截面.

例如，设 $\frak{su}(N)$ 表示 $SU(N)$ 的李代数. $\frak{su}(2)$ 到 $\frak{su}(N)$ 的 Kostant 嵌入，是使 $\frak{su}(N)$ 的 $N$ 维表示在 $\frak{su}(2)$ 下不可约的嵌入（因此它按自旋 $(N-1)/2$ 变换）. 在群的层面上，对于 $SU(N)$，当 $N$ 为奇数时，Kostant 映射的像为 $SO(3)$；当 $N$ 为偶数时，则为 $SU(2)$. 更一般地，如果 $G=SU(N)/\Gamma$，其中 $\Gamma$ 是 $SU(N)$ 的中心的子群，那么除非 $N$ 为偶数且 $G$ 的元素 $-1$ 不属于 $\Gamma$，Kostant 映射的像都为 $SO(3)$. 例如，对于 $G=PSU(N)$，Kostant 映射的像始终为 $SO(3)$. 这些结果当然与我们在第 \ref{whichline} 节中的发现一致.

一般而言，如果 $G$ 是伴随型的，那么 Kostant 映射的像始终为 $SO(3)$. 所以，例如对于 $G=E_8$，希钦纤维化的纤维典范地是一个阿贝尔簇.





\section{对偶环面与希钦纤维化}
\label{dualtori}
\markright{\thesection. 对偶环面与希钦纤维化}

正如我们在第 \ref{compacthitchin} 节中所回顾的，对于任意单李群，由于空间 $\MH(G,C)$ 完全可积，希钦纤维化\index{hitchin-fibration@希钦纤维化}的纤维 $\FF$ 一般是一个复环面. 物理上的论证 \cite{BJSV,HMS} 表明，四维中的 $S$-对偶性通过希钦纤维化纤维 $\FF$ 上的 $T$-对偶性以及底空间 ${\cal V}$ 的一个几何对称性\footnote{对于我们作为例子讨论的酉群，底空间的几何对称性只是一个简单的缩放，不起重要作用.}起作用. 正如 \cite{HauselThaddeus} 最初解释的，$\MH(^L\neg G,C)$ 与 $\MH(G,C)$ 之间的镜像对称\index{mirror-symmetry@镜像对称}是能够相当详细地实施镜像对称的 SYZ 方法 \cite{SYZ} 的一个罕见例子，因为超凯勒\index{hyperkahler@超凯勒}结构使我们能够相当显式地构造以拉格朗日环面为纤维的纤维化. 对 $\MH(^L\neg G,C)$ 与 $\MH(G,C)$ 之间对偶性的这一解释，在 \cite{KW} 中发挥了重要作用. 这里我们将对其中一些方面给出比该文更详尽的解释.


确切地说，我们的断言是：若 $G$ 和 $^LG$ 是两个对偶群，那么位于底空间中相应点之上的希钦纤维 $\FF$ 和 $^L\neg \FF$ 是对偶复环面. 回顾一下，这意味着 $\FF$ 参数化 $^L\neg \FF$ 上的平坦线丛\index{flat-line-bundle@平坦线丛}，反之亦然. 描述复环面 $\cF$ 与 $\cF'$ 之间的对偶性最对称的方式，是给出乘积 $\cF\times
{}\cF'$ 上一个带联络的酉线丛\index{line-bundle@线丛} ${\cal T}$，使得对任意 $f\in \cF$，${\cal T}$ 限制到 $f\times
{}\cF'$ 时是平坦的，并且对任意 $f'\in \cF'$，限制到 $\cF\times {}f'$ 时也是平坦的.（如果 $\cF$ 和 $\cF'$ 具有复结构\index{complex-structure@复结构}或代数结构，则要求 ${\cal T}$ 是全纯线丛或代数线丛.）因此，让 $f$ 变化，就得到 $\cF'$ 上的一族平坦线丛 ${\cal Y}'$；让 $f'$ 变化，就得到 $\cF$ 上的一族平坦线丛 ${\cal Y}$. 如果 ${\cal Y}'$ 是 $\cF'$ 上平坦丛的模空间\index{moduli-space@模空间}，且 ${\cal Y}$ 是 $\cF$ 上平坦丛的模空间，那么称 $\cF$ 和 $\cF'$ 为对偶环面\index{dual-torus@对偶环面}，并称 ${\cal T}$ 为这对环面的一个庞加莱线丛\index{poincare-line-bundle@庞加莱线丛}.

该定义的一个直接推论是，如果 $\cF$ 和 $\cF'$ 是对偶环面，那么总存在唯一的特选点 $f\in \cF$，使得 ${\cal T}$ 限制到 $f\times {}\cF'$ 时平凡；反之，总存在唯一的特选点 ${}f'\in{} \cF'$，使得 ${\cal T}$ 限制到 $\cF\times {}f'$ 时平凡. 因此，如果 $\cF$ 和 $\cF'$ 是在上述意义下对偶的复代数环面，它们总是典范地成为阿贝尔簇\index{abelian-variety@阿贝尔簇}.\transnote{“唯一特选点”给出一种标准的主极化：把 $\cF$ 与 $\cF'$ 的格分别看作电荷格与磁荷格，固定 ${\cal T}$ 在唯一特选点处的平凡化，便可把 $\cF'$ 等同于 $\mathrm{Pic}^0(\cF)$。这正是 Goddard–Nuyts–Olive 对偶（GNO 对偶）的几何形式：一个群的电荷格与对偶群的磁荷格互为对偶。（来源：A 级，Goddard–Nuyts–Olive，\emph{Gauge Theories and Magnetic Charge}，Nucl. Phys. B125 (1977) 1–28；原文 §3 对偶环面定义；2026-11-19 核对.）}

本节致力于在希钦纤维化的背景下讨论这些问题. 对于其中若干点，我们将给出比本文其余部分所需更详细的讨论.

\subsection{例子}
\label{examplo}

作为对偶环面的一个初等例子，设 $S$ 是由角变量 $x$、$0\leq x\leq 2\pi$ 参数化的圆，$S'$ 是由角变量 $y$、$0\leq y\leq 2\pi$ 参数化的另一个圆. 当然，$S\times S'\cong
\Bbb{R}^2/\Bbb{Z}^2$，其中 ${\Bbb{Z}}^2$ 通过 $(x,y)\to (x+2\pi
n,y+2\pi m)$ 作用. 令 $\omega= dx\wedge dy/2\pi$. 曲率为 $F=-i\omega$ 的线丛 ${\cal T}\to S\times S'$ 可以充当庞加莱线丛. 我们可以如下定义这样一个 ${\cal T}$：在 $\Bbb{R}^2$ 上的平凡线丛 ${\cal O}$ 上选择联络形式 $B=- ix\,dy/2\pi$，然后除以 $x$ 和 $y$ 方向的单位平移，将其下降到 $S\times S'$.（我们令 $(x,y)\to (x+2\pi n,y+2\pi m)$ 在 ${\cal O}$ 上的作用为乘以 $e^{ iny}$.）${\cal T}$ 绕圆 $x\times S'$ 的和乐\index{holonomy@和乐}为 $e^{ ix}$，绕圆 $S\times y$ 的和乐为 $e^{iy}$. 平坦线丛的每一种可能和乐都在每个族中恰好出现一次，因此 ${\cal T}$ 是一个庞加莱线丛，它建立了 $S$ 与 $S'$ 之间的对偶性.

为了构造线丛 ${\cal T}$，除了选取一个显式的联络形式 $B$，我们也可以仅要求 ${\cal T}$ 的曲率为 $F=-i\omega$，并且限制到 $x=0$ 或 $y=0$ 时和乐平凡，从而达到同样的目的. 这些条件在同构意义下完全刻画了 ${\cal T}$.

这个初等例子有如下推广. 设 $\Gamma$ 和 $\Gamma^*$ 是任意两个对偶格，带有一个幺模\index{unimodularity@幺模}双线性形式 $\langle ~,~\rangle:\Gamma\times
\Gamma^*\to \Bbb{Z}$. 定义向量空间 $V=\Gamma{\otimes}_{\Bbb{Z}}\Bbb{R}$、$V^*=\Gamma^*{\otimes}
_{\Bbb{Z}}\Bbb{R}$ 以及环面 $\cF=V/\Gamma$、$\cF^*=V^*/\Gamma^*$. $V$ 和 $V^*$ 是对偶向量空间，因为形式 $\langle~,~\rangle$ 自然地延拓为双线性配对 $V\otimes V^*\to \Bbb{R}$. 此外，环面 $\cF$ 和 $\cF^*$ 总是以自然方式成为对偶环面. 只需从辛形式 $\omega=\langle dx,dy\rangle/2\pi$ 出发，进行我们刚才见过的相同构造即可，这里我们用 $x$ 或 $y$ 表示 $\cF$ 或 $\cF^*$ 中的点到 $V$ 或 $V^*$ 的提升.

举一个更接近我们主题的例子，设 $D$ 是紧致的二维定向闭流形——暂时不赋予复结构——并设 ${\cal J}$ 是 $D$ 上平坦酉复线丛的模空间. 那么 ${\cal J}$ 是自对偶环面，因为它等于 $(\Gamma\otimes_{\Bbb{Z}}\Bbb{R})/\Gamma$，其中 $\Gamma$ 是自对偶格 $\Gamma= H_1(D,\Bbb{Z})$. 因而，可以应用上一段的讨论，得到 ${\cal J}$ 的一个典范自对偶性.

此外，为了有助于后面的工作，我们也可以用规范理论\index{gauge-theory@规范理论}以下述方式理解 ${\cal J}$ 的自对偶性. 我们如下定义线丛 ${\cal T}\to {\cal J}\times {\cal J}$. 设 ${\cal L}$ 是 $D$ 上带联络 $A$ 的平坦线丛，并设 ${\cal M}$ 是 $D$ 上带联络 $B$ 的第二个平坦线丛. 由于平坦，这些联络满足
\begin{equation}
\label{olpp}
0=F_A=dA,~~0=F_B=dB.
\end{equation}
这对 ${\cal L}$、${\cal M}$ 确定一点 $P\times P'\subset
{\cal J}\times {\cal J}$. 我们通过如下公式定义 ${\cal J}\times {\cal J}$ 上的一个二形式
\begin{equation}
\label{ivo} \omega={1\over 2\pi}\int_{D}\delta A\wedge \delta B.
\end{equation}
若用一个生成 $A$ 或 $B$ 的规范变换的向量场与 $\omega$ 缩并，结果为零，因为若在 $\delta A$ 或 $\delta B$ 上加一个恰当形式 $d\epsilon$ 并分部积分，利用平坦性条件 (\ref{olpp}) 就得到零. 因此，$\omega$ 下降为 ${\cal J}\times {\cal J}$ 上的二形式，而由庞加莱对偶性，这个形式是非退化的. 对任意 $P,P'\subset {\cal J}$，$\omega$ 限制到 $P\times {\cal J}$ 时为零（这可通过令 $\delta A=0$ 实现），限制到 ${\cal J}\times P'$ 时也为零（通过令 $\delta B=0$ 实现）. 此外，我们已将 $\omega$ 归一化，使它的周期为 $2\pi$ 的整数倍.

由这些性质可知，存在 ${\cal J}\times {\cal J}$ 上带联络的酉线丛 ${\cal T}$，使得 ${\cal T}$ 的曲率等于 $-i\omega$. 此外，如果要求 ${\cal T}$ 限制到 $0\times
{\cal J}$ 或 ${\cal J}\times 0$ 时平凡（平坦且单值变换\index{monodromy@单值变换}为零），那么它在同构意义下唯一确定，其中 $0\in {\cal J}$ 是对应于平凡平坦联络\index{flat-connection@平坦联络}的点.

这样的 ${\cal T}$ 建立了 ${\cal J}$ 的自对偶性，即 ${\cal J}$ 是它自身的对偶环面. 事实上，如果在 $D$ 上选取一维闭链的辛基，使交叉形式化为形如 $\begin{pmatrix}0
& 1\\ -1 & 0\end{pmatrix}$ 的 $2\times 2$ 块之和，那么这个例子就化为我们最初考虑的初等例子的若干副本的乘积，即 $S\times S'$ 上曲率为 $-idx\wedge dy/2\pi$ 的线丛.

这也使我们能够用初等公式显式构造 ${\cal T}$（就像在 $S\times S'$ 的简单模型例子中所做的那样）. 不过，构造 ${\cal T}$ 的一种更优雅的方式，是考虑三维流形 $\Bbb{R}\times D$ 上的陈-西蒙斯规范理论，其中规范群为 $U(1)\times U(1)$，规范场\index{gauge-field@规范场}由一对 $U(1)$ 规范场组成，我们记它们为 $A,B$. 取作用量为
\begin{equation}
\label{ploko} I={1\over 2\pi}\int_{\Bbb{R}\times D}\,A\wedge dB.
\end{equation}
经典运动方程告诉我们 $dA=dB=0$，而经典相空间是解模去规范变换后的模空间，即 ${\cal J}\times {\cal J}$. 经典相空间上的辛结构就是我们所记的 $\omega$，而量子化的第一步是构造具有我们所宣称性质的线丛 ${\cal T}$. 这可以利用陈-西蒙斯作用量的规范不变性直接完成（例如，见 \cite{rsw,axelrod}）.\transnote{其机制是：陈-西蒙斯作用量 $I=(1/2\pi)\int A\wedge dB$ 在带边三维流形上的变分给出边界项 $\frac{1}{2\pi}\int_{\partial X}\delta A\wedge\delta B$，故 $\omega$ 正是 $I$ 的“边界辛形式”；$1/2\pi$ 的归一化使 $I$ 的周期属于 $2\pi\Bbb Z$，从而预量子化线丛 ${\cal T}$ 存在且在同构意义下唯一。（来源：A 级，Ramadas–Singer–Weitsman，\emph{Some Comments On Chern-Simons Gauge Theory}，Commun. Math. Phys. 126 (1989) 409；原文 §3.1 式 (3.1)–(3.3)；2026-11-19 核对.）}

与许多可以在一般光滑二维曲面 $D$ 上进行的构造一样，如果在 $D$ 上选定一个复结构，这一构造也可以有一种重要的其他描述. 此时，${\cal J}$ 成为 $D$ 的雅可比簇\index{jacobian@雅可比簇}，它参数化 $D$ 上次数为零的全纯线丛；我们将其记为 ${\rm Jac}(D)$. 这是自对偶阿贝尔簇的一个经典例子，并且与希钦纤维化密切相关，我们很快就会讨论这一点. 除了上述拓扑构造，${\rm Jac}(D)$ 的庞加莱线丛也可以方便地用全纯方法构造；我们将在第 \ref{ugduality} 节解释其中一种做法.

\subsection{酉群的情形} \label{unitaryexample}

这里我们将显式描述对偶群对 $SU(N)$ 与 $PSU(N)$ 所对应的对偶希钦纤维化的纤维之间的对偶性，以及类似地，$U(N)$ 的自对偶性.

这样做时，我们立即遇到如下事实：$\MH(PSU(N))$ 有 $N$ 个连通分支，以一个离散特征类\index{characteristic-class@特征类} $\xi$（在第 \ref{psun} 节中引入）的取值标记，而 $\MH(SU(N))$ 是连通的. 正如 \cite{HauselThaddeus} 以及 \cite{KW} 第 7 节所解释的（另见下文第 \ref{gerbes} 节），$\xi$ 的对偶是可以在 $\MH(SU(N))$ 上引入的某个离散 $B$-场（或 gerbe）. 我们将把这个问题留到后面，暂时只考虑 $\MH(SU(N))$ 与 $\MH(PSU(N))$ 中 $\xi=0$ 的连通分支 $\MH^{(0)}(PSU(N))$ 之间的对偶性. 类似地，在 $U(N)$ 的情形下，我们暂时只考虑第一陈类\index{chern-class@第一陈类}为零的希格斯丛\index{higgs-bundle@希格斯丛}.（第 \ref{dualap} 节给出不带这些限制的对偶性表述.）



即使有了这些限制，还会出现另一个难题. 正如上面解释的，如果 $\cF$ 和 $\cF'$ 是对偶复代数环面，它们总是典范地成为阿贝尔簇. 纤维 $\FF^{(0)}_{PSU(N)}$ 对所有 $N$ 都有这一性质；但正如我们已经看到的，$\FF_{SU(N)}$ 和 $\FF_{U(N)}$ 只有在 $N$ 为奇数时才自然地成为阿贝尔簇. 对于偶数 $N$，只有在 $C$ 上选定自旋结构\index{spin-structure@自旋结构} $K^{1/2}$ 后，它们才自然地成为阿贝尔簇. 因而，$\FF_{SU(N)}$ 与 $\FF^{(0)}_{PSU(N)}$ 对偶的断言（以及类似地，$\FF_{U(N)}$ 自对偶的断言），对于偶数 $N$ 不能算是完全正确的. 文献 \cite{FW} 第 10 节给出了修正后的表述，但这里不作介绍. 就我们的目的而言，只需说，对于偶数 $N$，一旦选定 $K^{1/2}$，关于 $\FF_{SU(N)}$、$\FF^{(0)}_{PSU(N)}$ 和 $\FF_{U(N)}$ 的对偶性断言就成立.\footnote{\label{brief} 简要地说，\cite{FW} 中的表述是：对于偶数 $N$，$S$-对偶性在 $\MH$ 上生成一个平凡但非典范平凡的 $B$-场（或 gerbe）；通过选择 $K^{1/2}$，可以将其平凡化\index{trivialization@平凡化}.（这个 gerbe 在第 \ref{gerbes} 节中描述.）对于偶数 $N$，若要在不选择 $K^{1/2}$ 的情况下表述 $\FF_{SU(N)}$ 与 $\FF^{(0)}_{PSU(N)}$ 之间的对偶性，或 $\FF_{U(N)}$ 的自对偶性，就必须在对偶性的一侧纳入这个 $B$-场. 关于为何 $S$-对偶性会生成这个平凡但非典范平凡的 $B$-场，目前尚未给出直接来自四维规范理论的解释.}


回顾一下，纤维 $\FF_{SU(N)}$ 由谱覆盖\index{spectral-cover@谱覆盖} $\psi:D\to C$ 确定，并且正如在第 \ref{whichline} 节中发现的，由次数为 $N(N-1)(g_C-1)$ 的线丛 $\L\to D$ 的选择来参数化. 一旦选定 $K^{1/2}$，$\FF_{SU(N)}$ 就包含典范点 $\psi^*(K^{(N-1)/2})$，从而成为阿贝尔簇，可以方便地用次数为零的线丛
\begin{equation}\label{zumbo}
{\cal L}'={\cal L}\otimes \zeta^*(K^{-(N-1)/2})\end{equation}
来参数化. 为了用 ${\cal L}'$ 描述 $\FF_{SU(N)}$，引入另一个概念会很有用. 给定黎曼曲面之间的覆盖 $\psi:D\to C$，可以定义一个“范”映射 ${\rm Nm}:{\rm Pic}(D)\to{\rm Pic}(C)$. 如果 ${\cal M}\to D$ 的除子\index{divisor@除子}为 $\sum n_i\qq_i$，其中 $n_i$ 为整数且 $\qq_i\in D$，那么 ${\rm Nm}({\cal M})$ 的除子为 $\sum_in_i\psi(\qq_i)$. 范映射\index{norm-map@范映射}与之前考虑的推前\index{pushforward@推前}操作 $\cal M\to \det\psi_*(\cal M)$ 的区别在于，它{
\it 没有}分歧修正项. 对于我们的 $N$ 重覆盖 $\psi:D\to C$，两者通过 $\det\psi_*({\cal M})={\rm Nm}({\cal M})\otimes
K^{-N(N-1)/2}$ 联系起来，对任意 ${\cal M}$ 都成立.（只需对 ${\cal M}={\cal O}$ 验证这一点；我们有 ${\rm Nm}({\cal O})={\cal O}$，而 (\ref{jofo}) 意味着 $\det\psi_*({\cal O})=K^{-N(N-1)/2}$.）因此，关于 $SU(N)$ 的希钦纤维化的纤维的描述 (\ref{yudo}) 等价于说：用 $\cal L'$ 参数化时，$\FF_{SU(N)}$ 是由下式定义的 ${\rm Jac}(D)$ 的子群
\begin{equation}\label{nirk}  \FF_{SU(N)}={\rm Nm}^{-1}({\cal O}).\end{equation}

另一方面，用 ${\cal L}'$ 表示时，拓扑平凡的 $PSU(N)$ 丛的希钦纤维化的纤维 $\FF^{(0)}_{PSU(N)}$ 为
\begin{equation}\label{lirk}
\FF^{(0)}_{PSU(N)}={\rm Jac}(D)/{\rm Jac}(C).\end{equation}
文献 \cite{HauselThaddeus} 以相当初等的方式证明了这两者对偶.

粗略的解释如下. 如果 ${\cal J}$ 是任意自对偶环面，带有庞加莱线丛 ${\cal P}\to {\cal J}\times
{\cal J}$，并且 $\cF$ 是 ${\cal J}$ 的任意子环面，那么 $\cF$ 的对偶为 ${\cal J}/\cF^\perp$，其中 $\cF^\perp$ 由所有满足如下条件的 $f^\perp\in {\cal J}$ 组成：${\cal P}$ 限制到 $\cF\times
f^\perp$ 时平凡. 对于 ${\cal J}={\rm Jac}(D)$、$\cF=\FF_{SU(N)}={\rm Nm}^{-1}({\cal O})$，有 $\cF^\perp=\psi^*({\rm Jac}(C))$，从而 $\cF$ 的对偶为 $\FF_{PSU(N)}^{(0)}={\cal J}/\cF^\perp={\rm Jac}(D)/{\rm Jac}(C)$. 我们将在第 \ref{topview} 节作更多解释，并在第 \ref{dualap} 节描述处理这种对偶性的另一种方法.

类似地，用 ${\cal L}'$ 代替 ${\cal L}$ 表示时，(\ref{blonko}) 变为
\begin{equation}\label{pirk}\FF^{(d)}_{PSU(N)}={\rm Pic}_d(D)/{\rm
Jac}(C).\end{equation}
通过重新表述 (\ref{normat})，可以得到一个等价的刻画：
\begin{equation}\label{dormat} \FF_{PSU(N)}^{(d)}={\rm
Nm}^{-1}({\cal L}_0)/\Gamma_N,\end{equation} 其中 ${\cal L}_0$ 是 $C$ 上任意固定的次数为 $d$ 的线丛. 同样，(\ref{palooko}) 变为
\begin{equation}\label{kormat}\hat \FF_{PSU(N)}^{(d)}={\rm
Nm}^{-1}({\cal L}_0).\end{equation}

对于 $U(N)$，类似的结论很简单：当用次数为零的线丛 ${\cal L}'$ 参数化时，希钦纤维化的纤维是 $D$ 的雅可比簇的一个副本：
\begin{equation}\label{zormo}\FF_{U(N)}={\rm Jac}(D).\end{equation}
这当然是一个自对偶阿贝尔簇.



\subsection{拓扑观点}
\label{topview}

这里，对于 $G=SU(N)$，我们从拓扑观点更仔细地考察谱曲线\index{spectral-curve@谱曲线}. 我们将补上前面论证中的一些缺口，并得到若干有用的结果. 此外，鉴于几何朗兰兹纲领可以由拓扑场论推导出来这一基本断言，我们希望尽可能从拓扑角度理解这些内容.

一般地，从一个黎曼曲面到另一个黎曼曲面的映射 $\psi:D\to C$ 确定一个相应的映射 $\psi_*:H_1(D,\Bbb{Z})\to
H_1(C,\Bbb{Z})$. 只需把闭路 $\gamma\subset D$ 的同调类映到相应闭路 $\psi(\gamma)\subset C$ 的同调类. 我们把同调群简记为 $H_1(D)$ 和 $H_1(C)$，并将此映射记为
\begin{equation}\label{orgok}\psi_*:H_1(D)\to
H_1(C).\end{equation}

如果 $C$ 中的每条闭路都可以提升为 $D$ 中的闭路，那么映射 $\psi_*$ 就是满射. 这并不是对黎曼曲面之间的每个映射都成立；例如，对无分歧\index{unramified@无分歧}覆盖 $D\to C$ 就不成立. 然而，谱覆盖 $\psi:D\to C$ 有大量分歧，在这种情况下 $\psi_*$ 是满射. 一种证明方法是选取一个方便的谱覆盖，利用光滑谱覆盖的空间是连通的这一事实，因此任意两个光滑谱覆盖都具有相同的拓扑. 我们选取一点 $\pp\in C$，再选取一个谱覆盖，使得 $\Tr\,\varphi^n$ 对 $n=2,\dots, N$ 都在 $\pp$ 处为零，并且 $\Tr\,\varphi^N$ 在 $\pp$ 处仅有一个单零点. 此时，谱覆盖的方程在 $\pp$ 附近形如 $y^N-z=0$，其中 $z$ 是 $C$ 上在 $\pp$ 处为零的局部参数. 因而，在这个例子中，$D$ 只包含唯一一个位于 $\pp$ 上方的点 $\qq$. 现在，$C$ 中的任意闭路都可以形变为一条以 $\pp$ 为起点和终点的路径. 这样的路径在 $D$ 中提升为一条以 $\qq$ 为起点和终点的路径，因此构成闭路. 所以映射 $\psi_*$ 是满射.

我们用 $\Gamma$ 表示 $\psi_*$ 的核. 由于 $\psi_*$ 是满射，我们得到格的正合列\index{exact-sequence@正合列}
\begin{equation}\label{helfgo}0\to \Gamma\to H_1(D)\underarrow{\psi_*}H_1(C)\to
0.\end{equation} 这样的正合列总可以分裂，尽管分裂并非典范：
\begin{equation}\label{norgo}H_1(D)=\Gamma\oplus
H_1(C).\end{equation} 相对于这一分裂，映射 $\psi_*$ 就是
\begin{equation}\label{orgo}\psi_*=0\oplus 1.\end{equation}

我们还可以定义另一个映射 $\psi^*:H_1(C)\to H_1(D)$，它把闭路 $\gamma\subset C$ 的同调类映到 $\psi^{-1}(\gamma)\subset D$ 的同调类. 对于任意 $\gamma\subset C$，$\psi(\psi^{-1}(\gamma))$ 是 $\gamma$ 的一个 $N$ 重覆盖，这在同调层面上意味着
\begin{equation}\label{ilbo}\psi_*\circ\psi^*=N.\end{equation}
因此，相对于分解 (\ref{norgo})，有
\begin{equation}\label{milbo}\psi^*=m\oplus N,\end{equation}
其中 $m:H_1(C)\to\Gamma$ 是某个格映射，而 $N:H_1(C)\to
H_1(C)$ 表示乘以 $N$. 由于分裂 (\ref{norgo}) 并非典范，$m$ 仅在模 $N$ 意义下唯一确定.\transnote{此处的 $\psi^*$ 与 §2.6.6、§3.2 中的 $\psi^*$ 同名而含义不同：(3.15) 的 $\psi^*:H_1(C)\to H_1(D)$ 把 $C$ 中的闭链提升为 $D$ 中的闭链，满足 $\psi_*\psi^*=N$（$N$ 为覆盖次数）；而线丛情形下的 $\psi^*:\mathrm{Pic}(C)\to\mathrm{Pic}(D)$ 是拉回。读数时须按上下文分辨。（来源：A 级，原文 §3.3 式 (3.15)(3.16) 与 §2.6.6 式 (2.36)–(2.39)；2026-11-19 核对.）}

正合列 (\ref{helfgo}) 的对偶是另一个正合列
\begin{equation}\label{plygo}0\to H_1(C)^*\underarrow{i} H_1(D)^*\to
\Gamma^*\to 0.\end{equation} 这里 $\Gamma^*={\rm
Hom}(\Gamma,\Bbb{Z})$ 是格 $\Gamma$ 的对偶格. 按定义，映射 $i$ 将 $x:H_1(C)\to \Bbb{Z}$ 映到 $x\circ
\psi_*:H_1(D)\to \Bbb{Z}$. 然而，$H_1(C)$ 和 $H_1(D)$ 通过各自的交叉配对\index{intersection-pairing@交叉配对}自对偶，我们将这些配对记为 $\langle~,~\rangle_C$ 和 $\langle ~,~\rangle_D$. 若对于某个 $x'$，$x:H_1(C)\to
\Bbb{Z}$ 为 $y\to \langle x',y\rangle_C$，则 $i(x)$ 将 $z\in H_1(D)$ 映到 $\langle x',\psi_*(z)\rangle_C$. 但对于任意 $x'\in H_1(C)$、$z\in H_1(D)$，我们有 $\langle
x',\psi_*(z)\rangle_C =\langle \psi^*(x'),z\rangle_D$（这可通过计算 $C$ 和 $D$ 上的交点来论证）. 这意味着，在 $H_1(D)$ 的自对偶性下，元素 $i(x)\in
H_1(D)^*$ 对应于 $\psi^*(x')\in H_1(D)$. 因此，我们可以把正合列 (\ref{plygo}) 写成一个更有用的形式：
\begin{equation}\label{lygo}
0\to H_1(C)\underarrow{\psi^*} H_1(D)\to \Gamma^*\to
0.\end{equation} 这告诉我们 $H_1(D)/\psi^* (H_1(C))$ 是一个格，事实上
\begin{equation}\label{polyfo}
H_1(D)/\psi^*(H_1(C))=\Gamma^*.\end{equation}

$C$ 上平坦 $U(1)$ 丛的模空间为 ${\cal J}(C)={\rm Hom}(H_1(C),U(1))$，我们将其简记为 $H_1(C)^\vee$. 同样，$D$ 上平坦 $U(1)$ 丛的模空间为 ${\cal J}(D)={\rm Hom}(H_1(D),U(1))=H_1(D)^\vee$. 存在自然映射 $(\psi_*)^\vee:H_1(C)^\vee\to H_1(D)^\vee$，它将 $\phi:H_1(C)\to U(1)$ 映到 $\phi\circ\psi_*:H_1(D)\to U(1)$. 拓扑平凡的 $PSU(N)$ 丛的希钦纤维化的纤维为
\begin{equation}\label{doorstop}
\FF^{(0)}_{PSU(N)}={\cal J}(D)/{\cal J}(C)=H_1(D)^\vee/(\psi_*)^\vee
(H_1(C)^\vee).\end{equation} 在 (\ref{norgo}) 的分裂下，$\FF^{(0)}_{PSU(N)}$ 参数化从 $H_1(D)$ 到 $U(1)$ 的同态，只是我们不关心这样的同态如何作用于 $H_1(C)$. 因而
\begin{equation}\label{hormstop}
\FF^{(0)}_{PSU(N)}={\rm Hom}(\Gamma,U(1))=\Gamma^\vee.\end{equation}

类似地，定义 $\psi^*{}^\vee:H_1(D)^\vee\to H_1(C)^\vee$，它将 $\phi:H_1(D)\to U(1)$ 映到 $\phi\circ\psi^*:H_1(C)\to U(1)$. 这是代数几何中范映射的拓扑版本. 在代数几何中，对任意次数的线丛都可以定义范映射，但对于次数零，它等价于 $\psi^*{}^\vee$，后者的定义不使用复结构. 根据 (\ref{nirk})，$SU(N)$ 的希钦纤维化的纤维是范映射的核，即
\begin{equation}\label{kerno}\FF_{SU(N)}={\rm
ker}(\psi^*{}^\vee).\end{equation} 具体地说，${\rm ker}(\psi^*{}^\vee)$ 参数化在 $\psi^*(H_1(C))$ 上平凡的同态 $\phi:H_1(D)\to
U(1)$. 这些同态等同于从 $H_1(D)/\psi^*(H_1(C))=\Gamma^*$ 到 $U(1)$ 的同态，因此实际上有
\begin{equation}\label{berno} \FF_{SU(N)}={\rm
Hom}(\Gamma^*,U(1))=(\Gamma^*)^\vee.\end{equation} 特别地，正如第 \ref{whichline} 节中所预告的，$\FF_{SU(N)}$ 是连通的，并且是一个环面.\transnote{拓扑观点的要点是：$\FF_{SU(N)}=\mathrm{Hom}(\Gamma^*,U(1))$，只要 $\Gamma^*=H_1(D)/\psi^*H_1(C)$ 无挠，$\mathrm{Hom}(\Gamma^*,U(1))$ 就是紧环面，特别地连通；“无挠”来自 $\psi_*$ 满射，它保证 $\psi^*H_1(C)$ 是 $H_1(D)$ 的饱和子格。这与 §2.6.3 由复几何得到的预言互为独立证明。（来源：A 级，原文 §3.3 式 (3.19)–(3.23)；A 级，Hausel–Thaddeus，arXiv:math/0205236 对阿贝尔簇纤维的拓扑处理；2026-11-19 核对.）}

比较 (\ref{hormstop}) 和 (\ref{berno}) 可见，$\FF^{(0)}_{PSU(N)} $ 和 $\FF_{SU(N)}$ 由对偶格 $\Gamma$ 和 $\Gamma^*$ 分别给出，即 ${\rm Hom}(\Gamma,U(1))$ 与 ${\rm Hom}(\Gamma^*,U(1))$. 正如第 \ref{examplo} 节中所解释的，以这种方式得到的环面总是典范对偶的.




由于 $\Gamma$ 是 $H_1(D)$ 的子格，$H_1(D)$ 上的幺模交叉配对限制到 $\Gamma$ 上，成为一个整数值双线性形式. 这给出自然的单射 $\Gamma\to \Gamma^*$. 如果 $r_N$ 表示模 $N$ 约化，那么存在自然的满射 $r=r_N\circ \psi_*:H_1(D)\to
\Gamma_N$，其中 $\Gamma_N=H_1(C,\Bbb{Z}_N)$. 由 (\ref{ilbo})，$r_N$ 在 $\psi^*(H_1(C))$ 上为零，因此可以将它视为从 $\Gamma^*=H_1(D)/\psi^*(H_1(C))$ 到 $\Gamma_N$ 的满射. 它在 $\Gamma$ 上也为零，因为 $\psi_*$ 在其上为零，所以它实际上给出从 $\Gamma^*/\Gamma$ 到 $\Gamma_N$ 的满射. 最后，通过取判别式\index{discriminant@判别式}，可以证明这个映射是同构. 因而得到正合列
\begin{equation}\label{trulfo}0\to
\Gamma\to\Gamma^*\to\Gamma_N\to 0.\end{equation} 取到 $U(1)$ 的同态，并利用 $\FF_{PSU(N)}^{(0)}= \Gamma^\vee$ 以及 $\FF_{SU(N)}=(\Gamma^*)^\vee$，得到
\begin{equation}\label{mikolfo}\FF^{(0)}_{PSU(N)}=\FF_{SU(N)}/\Gamma_N.\end{equation}
这一结果已经在 (\ref{mormot}) 中从复几何的观点作了解释.

\subsubsection{$\FF_{SU(N)}$ 的刻画}
\label{character}

现在，我们回到复几何，看看从 $\FF_{SU(N)}={\rm Nm}^{-1}({\cal O})$ 是复环面这一事实可以推出什么.

设 $\pp$ 是 $C$ 中任意一点，而 $\qq,\qq'$ 是 $D$ 中位于其上方的两个点. 那么，$D$ 上的线丛 ${\cal O}(\qq)\otimes {\cal
O}(\qq')^{-1}$ 的范为 ${\cal O}$，因而定义 $\FF_{SU(N)}$ 中的一个点.

更一般地，设 $\pp_i$、$i-1,\dots,k$ 是 $C$ 中任意一组点，对每个 $i$，设 $\qq_i$ 和 $\qq'_i$ 位于 $D$ 中 $\pp_i$ 的上方. 选取整数 $n_i$ 并令 ${\DD}=\sum_in_i(\qq_i-\qq'_i)$. 那么线丛 ${\cal L}={\cal
O}({\DD})=\otimes_{i=1}^k\left({\cal O}(\qq_i)\otimes {\cal
O}(\qq'_i)^{-1}\right)^{n_i}$ 的范平凡，并属于 $\FF_{SU(N)}$.

我们断言，反过来，如果 ${\rm Nm}({\cal L})={\cal O}$，那么对于这种形式的某个除子 ${\DD}$，有 ${\cal L}={\cal O}({\DD})$. 论证依赖于已知 $\FF_{SU(N)}$ 是复环面. 既然如此，就可以除以 $N$，找到某个 ${\cal M}\in \FF_{SU(N)}$，使得 ${\cal M}^N\cong {\cal L}$. 存在某个除子 ${\DD}'=\sum_i n_i\qq_i$，使得 ${\cal M}={\cal O}({\DD}')$，因而 ${\cal L}={\cal O}(N{\DD}')$. 现在，对每个 $i$，设 $\qq_{i,\alpha}$、$\alpha=1,\dots, N$ 是 $D$ 中满足 $\psi(\qq_{i,\alpha})=\psi(\qq_i)$ 的点（其中一个 $q_{i,\alpha}$ 等于 $q_i$；为涵盖分歧，允许某些 $\qq_{i,\alpha}$ 相同）. 由于 $\cal M\in \FF_{SU(N)}$，${\rm Nm}({\cal M})$ 平凡，因此 $\psi^*({\rm Nm}({\cal M}))$ 也平凡. 但 $\psi^*({\rm Nm}({\cal M}))={\cal O}({\DD}'')$，其中 ${\DD}''=\sum_{i,\alpha}n_i \qq_{i,\alpha}$. 因此，$\psi^*({\rm Nm}({\cal M}))={\cal O}$ 这一事实意味着，可以等价地把 ${\cal L}$ 刻画为 ${\cal O}(N{\DD}'-{\DD}'')$. 但
\begin{equation}\label{helbo}N{\DD}'-{\DD}''=\sum_{i,\alpha}n_i(\qq_i-\qq_{i,\alpha}).\end{equation} 这是满足 $\psi(\qq_i)=\psi(\qq'_i)$ 的除子 $\qq_i-\qq_i'$ 之和，因此我们证明了每个 ${\cal L}\in {\rm Nm}^{-1}({\cal O})$ 都等于这种形式的某个除子 ${\DD}$ 所对应的 ${\cal O}({\DD})$.

\subsection{辛形式}
\label{sympform}

我们刚刚得到的结果有助于理解复结构 $I$ 下 $\MH$ 的几何的另一个方面. 至少对于酉群，我们希望用谱覆盖描述全纯辛结构\index{holomorphic-symplectic-structure@全纯辛结构} $\Omega_I$，它是在式 (\ref{zolk}) 中通过规范理论引入的.（这一详细描述不会在本文其余部分中使用.）对于更一般的群，人们希望能基于一种更抽象的谱覆盖方法得到类似结果.

我们的方法分两步进行. 首先，在由谱覆盖 $\psi:D\to C$ 和适当次数的线丛 ${\cal L}\to D$ 组成的对的空间上，直接构造一个全纯辛结构. 然后，将这个公式与规范理论定义所带来的预期进行比较.

首先回顾阿贝尔–雅可比映射\index{abel-jacobi-map@阿贝尔–雅可比映射}，它将黎曼曲面 $D$ 的雅可比簇显式地等同于一个复环面. 设 $g_D$ 是 $D$ 的亏格，并选取由 $g_D$ 个全纯 1-形式 $\omega^a$、$a=1,\dots, g_D$ 组成的基. 将表示 $H_1(D,\Bbb{Z})$ 中元素的整系数一维闭链 $\bar \gamma$ 映到周期 $(\int_{\bar\gamma}\omega^1,\dots,\int_{\bar
\gamma}\omega^{g_D})$，便把第一同调群 $H_1(D,\Bbb{Z})$ 映到 $\Bbb{C}^{g_D}$.（按定义，整系数一维闭链是定向闭路的整系数形式线性组合.）$H_1(D,\Bbb{Z})$ 在 $\Bbb{C}^{g_D}$ 中的像是秩为 $2g_D$ 的格 $\Gamma_D$. 商 ${\cal T}=\Bbb{C}^{g_D}/\Gamma_D$ 是一个复环面.

现在假设 ${\cal L}$ 是 $D$ 上次数为零的线丛. 它同构于某个除子 ${\DD}=\sum_in_i\qq_i$ 所对应的 ${\cal O}({\DD})$，其中整数 $n_i$ 满足 $\sum_in_i=0$，并且 $\qq_i\in D$. 由于 $\sum_in_i=0$，可以找到一个一维“链”$\gamma$（由不必闭合的定向路径构成的整系数形式线性组合），其边界 $\partial\gamma$ 等于 $\sum_i
n_i\qq_i$. 然后把 ${\cal L}$ 映到坐标为 $(\int_\gamma \omega^1,\dots,\int_\gamma
\omega^{g_D})$ 的点 $x\in {\cal T}$. 这给出 ${\cal T}$ 中一个良定义的点，因为 $\gamma$ 在相差一个一维闭链 $\bar\gamma$ 的意义下唯一确定，而将该闭链加到 $\gamma$ 上，会使 $x$ 的坐标平移格 $\Gamma_D$ 中的一个元素. 阿贝尔–雅可比定理表明，点 $x$ 只依赖于 ${\cal L}$，而不依赖于表示 ${\cal L}$ 的除子 ${\DD}$ 的选择，并且这个映射给出 ${\rm Jac}(D)$ 与复环面 ${\cal T}$ 的同构.

对于谱曲线，这可以用一种特别好的方式来实现. 我们从 $G=U(N)$ 的情形开始，再描述 $SU(N)$ 和 $PSU(N)$ 所需的微小变化. 在 $C$ 的余切丛\index{cotangent-bundle@余切丛}的全空间 $W$ 上，有一个自然的全纯 1-形式 $\lambda$，利用 $C$ 上的局部坐标 $z$，可以写成 $\lambda=y\, dz$. 一旦选定谱曲线 $D$，就可以将 $\lambda$ 限制为 $D$ 上的全纯微分. 当然，$D$ 由方程 $\det(y-\varphi)=0$ 定义，更显式地写成 $\PP(y,z;H_a)=0$，其中 $\PP$ 是关于 $y$ 的 $N$ 次多项式，而对易哈密顿量\index{hamiltonian@哈密顿量} $H_a$ 作为参数出现在其中.

$\lambda$ 限制到 $D$ 所得的上同调类依赖于参数 $H_a$. 只需在固定 $z$ 的情况下对 $\lambda$ 关于 $H_a$ 求导，就可以计算这种依赖. 以这种方式对方程 $\PP(y,z;H_a)=0$ 求导，得到 $0=\PP'(\partial y/\partial H_a)+\partial
\PP/\partial H_a$，其中 $\PP'$ 是 $\partial
\PP/\partial y$ 的简写. 因而
\begin{equation}\label{zonko} {\partial y\over \partial
H_a}=-{1\over \PP'}{\partial \PP\over \partial H_a}\end{equation}
所以，若令 $\omega^a=\partial \lambda/\partial H_a$，就得到
\begin{equation}\label{onko} \omega^a={\partial y\over \partial
H_a}dz=-{\partial \PP\over\partial H_a}{dz\over
\PP'}.\end{equation} 这些对象 $\omega^a$ 实际上是 $D$ 上的全纯 1-形式. 这里唯一不完全显然的一点是，$\omega^a$ 在 $\PP'=0$ 处没有极点. 但事实上，对于光滑平面曲线 $\PP(y,z)=0$，微分 $dz/(\partial \PP/\partial y)$ 在满足 $\PP=\partial \PP/\partial y=0$ 的点处是正则的，因为它也可以写成 $-dy/(\partial \PP/\partial z)$.

对于 $G=U(N)$，对易哈密顿量的数目等于 $D$ 的亏格 $g_D$. (\ref{onko}) 中定义的 $\omega^a$ 全部线性无关，因为多项式 $\PP(y,z;H_a)$ 确实非平凡地依赖于所有 $H_a$. 因而它们给出 $D$ 上全纯微分的一组基.

现在，我们几乎凭观察就可以写出 $\MH$ 上的一个全纯二形式. 若 $D$ 是谱曲线，${\cal L}$ 是 $D$ 上次数由 (\ref{ywo}) 给出的线丛，则写成 ${\cal L}=\psi^*(K^{(N-1)/2})\otimes {\cal L}'$，其中 ${\cal L}'$ 的次数为零. 我们用除子 ${\DD}=\sum_{i=1}^kn_i\qq_i$ 表示 ${\cal L}'$，并找到满足 $\partial\gamma={\DD}$ 的一维链 $\gamma$. 然后令
\begin{equation}\label{dunko}
\Omega=\sum_{a=1}^{g_D} dH_a\wedge d\int_\gamma
\omega^a.\end{equation} 为了看出这不依赖于 $\gamma$ 的选择，注意，若在 $\gamma$ 上加一个一维闭链 $\bar\gamma$，则 $\Omega$ 的变化为
\begin{align}\label{unko}
\Delta\Omega& =\sum_{a=1}^{g_D} dH_a\wedge d\int_{\bar\gamma}
\omega^a =\sum_{a,b=1}^{g_D}dH_a\wedge dH_b\int_{\bar
\gamma}{\partial \omega^a\over \partial H_b}\\
\nonumber &=\sum_{a,b=1}^{g_D}dH_a\wedge dH_b\int_{\bar
\gamma}{\partial^2 \lambda\over\partial H_a\partial
H_b}=0.\end{align} 类似地，定义 ${\cal L}'$ 时对 $K^{1/2}$ 的选择也无关紧要. 改变 $K^{1/2}$，其效果是用一个平方平凡的线丛来改变 ${\cal L}'$. 因此，它使周期向量 $(\int_\gamma\omega^1,\dots,\int_\gamma\omega^{g_D})$ 改变半个格向量，换言之，改变某个一维闭链 $\bar\gamma$ 所对应的 $(\int_{\bar
\gamma}\omega^1,\dots,\int_{\bar \gamma}\omega^{g_D})$ 的 $1/2$. 在 (\ref{unko}) 已进行的计算中再加一个因子 1/2，就说明这一操作不改变 $\Omega$.

实际上，我们可以用一种明显不依赖于一维链 $\gamma$ 选择的方式写出 $\Omega$. 在 (\ref{dunko}) 中，当外微分作用于 $\int_\gamma \omega^a$ 时，它可以对 $\gamma$ 求导，也可以对 $\omega^a$ 求导. 然而，按照与 (\ref{unko}) 相同的推理，对 1-形式 $\omega^a$ 求导的项不作贡献. 因此，只需对链 $\gamma$ 求导，或者更准确地说，对其由 $\partial\gamma=\sum_i
n_i\qq_i$ 刻画的端点求导. 因而有
\begin{equation}\label{wellwrite}\Omega=\sum_{a=1}^{g_D}\sum_{i=1}^kn_i
\,dH_a \wedge (\omega^a,d{\qq_i}).\end{equation} 符号 $d \qq_i$ 的含义如下. 由于 $\FF_{SU(N)}$ 的一个切向量可以具体地用 $\qq_i$ 的一阶位移表示，对于每个 $i=1,\dots,N$，存在从 $\FF_{SU(N)}$ 的切空间到 $D$ 在 $\qq_i$ 处的切空间 $TD_{\qq_i}$ 的自然映射. 等价地，对于每个 $i$，在 $\FF_{SU(N)}$ 上有一个取值于 $TD_{\qq_i}$ 的自然 1-形式. 这就是记为 $d \qq_i$ 的对象. 此外，$(\omega^a,\cdot)$ 表示 $D$ 上的 1-形式 $\omega^a$ 与 $TD$ 的配对. 因此，对每个 $i$，$(\omega^a,d \qq_i)$ 是 $\FF_{SU(N)}$ 上的一个 1-形式. 从而 $\Omega$ 如我们所需，是 $\M_H$ 上的 2-形式.



如果将希钦纤维化的纤维上的线性坐标记为 $X$，则 $\Omega$ 示意性地具有 $dX\wedge dH$ 的形式.（$X$ 的分量是 $q_i$ 的函数.）这保证函数 $H_a$ 在泊松意义下对易，并且通过泊松括号\index{poisson-bracket@泊松括号}生成各个 $X$ 的线性运动. 此外，在对希格斯场\index{higgs-field@希格斯场} $\varphi$ 进行缩放时，形式 $\Omega$ 是 $1$ 次齐次的（因为 $\lambda=y\,dz$ 具有这一性质，并且 $dH_a\,\omega^a=dH_a\,\partial\lambda/\partial H_a$ 与 $\lambda$ 具有相同的缩放行为）. 这些性质与底层规范理论中所定义的辛形式 $\Omega_I$ 的性质一致. 第 \ref{compgauge} 节将直接证明它们相同.



不过，先考虑 $G=SU(N)$ 和 $PSU(N)$ 的类似公式. 在亏格为 $g_C$ 的黎曼曲面 $C$ 上，有 $g_C$ 个全纯微分 $w^\alpha$、$\alpha=1,\dots,g_C$. 在 $D$ 上的 $g_D$ 个全纯微分中，可以取其中 $g_C$ 个为来自 $C$ 的拉回\index{pullback@拉回} $\omega^\alpha=\psi^*(w^\alpha)$. 还剩下 $D$ 上 $g_D-g_C$ 个并非来自 $C$ 的拉回的微分 $\tilde \omega^a$、$a=1,\dots,g_D-g_C$. 可以将 $\tilde\omega^a$ 归一化，要求对所有 $a$ 都有 $\psi_*(\tilde\omega^a)=0$. 这里，对于 $D$ 上的全纯微分 $\omega=f(y,z)\,dz$，$\psi_*(\omega)$ 是将 $\omega$ 推前到 $C$ 得到的. 具体地说，如果 $y_i$、$i=1,\dots,N$ 是方程 $\PP(y,z)=0$ 的根（将其视为固定 $z$ 后关于 $y$ 的方程），那么 $\psi_*(\omega)=\sum_i f(y_i,z)\,dz$.\transnote{$\psi_*$ 是有限态射 $\psi:D\to C$ 的迹映射（在每根纤维上取迹），它满射，核的维数恰为 $g_D-g_C$。对 $SU(N)$ 有 $\Tr\varphi=0$，故 $\lambda=y\,dz$ 满足 $\psi_*(\lambda)=\sum_i y_i\,dz=0$，即 $\lambda$ 落在核中；于是 $\partial\lambda/\partial H_a$（$a=1,\dots,g_D-g_C$）给出核的一组基。（来源：A 级，原文 §3.4 式 (3.32)(3.33)；A 级，Hitchin，\emph{Stable bundles and integrable systems}，Duke Math. J. 54 (1987) 91–114 关于谱曲线与迹映射；2026-11-19 核对.）}

从 $U(N)$ 规范理论转到 $SU(N)$ 规范理论时，通过令 $\Tr\,\varphi=0$ 去掉 $N$ 个对易哈密顿量. 特征多项式\index{characteristic-polynomial@特征多项式} $\PP(y,z;H_a)$ 的根 $y_i$ 与 $\varphi$ 的本征值相同，而 $\Tr\,\varphi=0$ 意味着 $\sum_i y_i=0$. 因此，全纯微分 $\lambda=y\,dz$ 满足 $\psi_*(\lambda)=\sum_iy_i\,dz=0$. 由于对于 $SU(N)$ 规范理论的对易哈密顿量的所有取值，都有 $\psi_*(\lambda)=0$，所以它对这些哈密顿量的导数也为零. 因而
\begin{equation}\label{milbox}0={\partial\over\partial
H_a}\psi_*(\lambda)=\psi_*\left({\partial\lambda\over
\partial H_a} \right),\end{equation}
这里 $H_a$ 是 $SU(N)$ 的对易哈密顿量，而不是 $U(N)$ 的. 因此，令 $\tilde\omega^a=\partial\lambda/\partial H_a$、$a=1,\dots, g_D-g_C$，得到的恰好是 $D$ 上被 $\psi_*$ 湮灭的那些微分.

现在将阿贝尔–雅可比映射应用于 $\FF_{SU(N)}$，它参数化满足 ${\rm Nm}({\cal L}')={\cal O}$ 的线丛 ${\cal L}'\in{\rm Jac}(D)$. 如第 \ref{character} 节所解释的，任何这样的 ${\cal L}'$ 都同构于 ${\cal O}(\DD)$，其中除子 ${\DD}$ 可取为 ${\DD}=\sum_in_i(\qq_i-\qq_i')$，满足 $\psi(\qq_i)=\psi(\qq_i')$. 对于满足 $\partial\gamma=D$ 的链 $\gamma$，我们可以取它为从 $\qq_i'$ 到 $\qq_i$ 的路径之和，从而 $\psi(\gamma)$ 是 $C$ 中的闭链. 如第 \ref{character} 节末尾所知，$C$ 中的这样的闭链可以提升为闭链 $\h\gamma\subset D$. 用 $\gamma-\h\gamma$ 替换 $\gamma$，便可以假定 $\psi(\gamma)=0$. 阿贝尔–雅可比映射将 ${\cal L}'$ 映到 $\Bbb{C}^{g_D}/\Gamma_D$ 中坐标为
\begin{equation}\label{gelmo}\left(\int_\gamma \tilde\omega^a,\int_\gamma
\psi^*(w^\alpha)\right).\end{equation} 的点. 然而，若 $w$ 是 $C$ 上的全纯微分，则 $\int_\gamma
\psi^*(w)=\int_{\psi(\gamma)}w=0$，因为我们要求了 $\psi(\gamma)=0$. 因此，${\cal L}'$ 在阿贝尔–雅可比映射下的像为
\begin{equation}\label{elmo}
\left(\int_\gamma\tilde\omega^a,~0\,\right).\end{equation} 将最后 $g_C$ 个坐标设为零后，现在只能用格 $\Gamma$ 中的元素来平移 $\gamma$，这个格是 $\psi_*$ 的核（回顾式 (\ref{helfgo})）. 用 $\Gamma$ 的元素平移 $\gamma$ 不影响阿贝尔–雅可比映射最后 $g_C$ 个分量为零，因为若 $\psi_*(\bar\gamma)=0$ 且 $w\in H^0(C,K)$，则 $\int_{\bar\gamma}\psi^*(w)=\int_{\psi_*(\bar\gamma)}w=0$. 用 $\Gamma$ 的任何其他元素平移 $\gamma$，则会破坏最后这 $g_C$ 个分量为零的性质.

因此，阿贝尔–雅可比映射将 $\FF_{SU(N)}$ 映到 $\Bbb{C}^{g_D-g_C}/\Gamma$. 这个映射是单射，因为即使在更大的空间 $\FF_{U(N)}$ 上，阿贝尔–雅可比映射也是单射；而从维数考虑，它也是满射. 所以，$\FF_{SU(N)}$ 同构于 $\Bbb{C}^{g_D-g_C}/\Gamma$. 这与 (\ref{berno}) 中 $\FF_{SU(N)}={\rm
Hom}(\Gamma^*,U(1))$ 的断言完全一致，因为 ${\rm Hom}(\Gamma^*,U(1))$ 典范地同构于 $(\Gamma\otimes_{\Bbb{Z}}\Bbb{R})/\Gamma$，而从 $C$ 的复几何观点来看，后者就是 $\Bbb{C}^{g_D-g_C}/\Gamma$.

至此，可以仿照 $U(N)$ 情形下辛形式的两个公式中的任意一个. 对于 $G=SU(N)$，通过下式定义 $\MH$ 上的全纯辛形式
\begin{equation}\label{flunko}
\Omega=\sum_{a=1}^{g_D-g_C} dH_a\wedge d\int_\gamma\tilde
\omega^a=\sum_{a=1}^{g_D-g_C}\sum_{i=1}^kn_i \,dH_a \wedge
(\tilde\omega^a,d\qq_i).\end{equation}

那么 $G=PSU(N)$ 又如何呢？$\MH(PSU(N))$ 的连通分支 $\MH^{(0)}(PSU(N))$ 参数化拓扑平凡情形下的 $PSU(N)$ 希格斯丛，它就是商 $\MH(SU(N))/\Gamma_N$，其中 $\Gamma_N$ 是 $C$ 上阶为 $N$ 的线丛所成的群. 用阶为 $N$ 的线丛扭曲 ${\cal L}'$ 时，辛形式 $\Omega$ 不变（在 (\ref{unko}) 的论证中加一个因子 $1/N$ 再重复一遍即可看出），所以相同的公式给出 $\MH^{(0)}(PSU(N))$ 上的全纯辛形式.

$\MH(PSU(N))$ 的其他连通分支同样是我们所称的 $\hat\MH^{(d)}(PSU(N))$（$\xi=d$ 时希钦方程\index{hitchin-equation@希钦方程}的解模去同伦于恒等变换的规范变换所得的模空间）除以 $\Gamma_N$ 的商，因此仍然只需在 $\hat\MH^{(d)}(PSU(N))$ 上构造辛形式. 选取一点 $\pp\in C$，令 ${\cal L}_0={\cal O}(\pp)^d$，则根据 (\ref{kormat})，希钦纤维化的纤维为 ${\rm Nm}^{-1}({\cal L}_0)$. 如果 ${\cal L}'$ 是满足 ${\rm Nm}({\cal L}')={\cal L}_0$ 的线丛，则对某个除子 ${\DD}'$ 写成 ${\cal L}'={\cal O}({\DD}')$，并令 ${\DD}=\DD'-d\qq$，其中 $\qq\in D$ 位于 $\pp\in C$ 上方. 因而 ${\rm Nm}(\O(\DD))=0$，所以除子 $\DD$ 线性等价于 $\DD''=\sum_in_i(\qq_i-\qq_i')$，其中 $\qq_i,\,\qq_i'\in D$ 在 $C$ 中具有相同的像. 找到满足 $\partial\gamma={\DD''}$ 的一维链 $\gamma$ 后，用相同的公式 (\ref{flunko}) 定义辛形式.

\subsubsection{与规范理论的比较}
\label{compgauge}

在规范理论中，全纯辛形式在式 (\ref{zolk}) 中定义为
\begin{equation}\label{zinto}\Omega_I={1\over\pi}\int_C|d^2z|\Tr\,\delta\phi_z
\delta A_{\bar z}.\end{equation} 回顾一下，符号 $\delta$ 表示无穷维函数空间（例如联络空间）上的外微分. 当降到有限维时，我们就用 $d$ 表示外微分.

元素 $\varphi_z \,dz\in H^0(C,\ad(E)\otimes K)$ 是将微分 $y\,dz\in H^0(D,K_D)$ 从谱覆盖 $D$ 向下推前得到的. 用 $D$ 上的数据，辛形式可以写成
\begin{equation}\label{minto}\Omega_I={1\over
\pi}\int_D|d^2z|\delta y\,\delta a_{\bar z},\end{equation} 其中线丛 ${\cal L}'\to D$ 的复结构由一个 $\bar\partial$ 算符定义，它在局部可以写成 $d\bar z(\partial_{\bar z}+a_{\bar z})$（在 $z$ 是良好局部参数的地方，即远离分歧点\index{ramification-point@分歧点}的地方）.

将它与 (\ref{wellwrite}) 这样的公式比较时，关键的一步是找到 $a_{\bar z}$ 的一种方便选择. 当然，$a_{\bar z}$ 只在模去规范变换的意义下唯一确定，而我们希望找到一个方便的代表元. 如果线丛 ${\cal L}'$ 是 ${\cal O}(\DD)$，其中 $\DD=\sum_in_i\rr_i$，$\rr_i\in D$，且 $\sum_in_i=0$，那么方程 $\bar\partial\partial f=i\pi\sum_in_i\,\delta_{\rr_i}$ 在 $D$ 上有模去加法常数后的唯一解. 这里 $\delta_{\rr_i}$ 是 $D$ 上的一个分布型 $(1,1)$-形式，支撑于 $\rr_i$，且满足 $\int_D\delta_{\rr_i}=1$. 为了简单起见，假设 $\rr_i$ 不是分歧点，从而 $z$ 是 $\rr_i$ 附近的良好局部参数，那么在 $\rr_i$ 附近有 $f\sim\frac {1}{
2}\ln|z-\rr_i|^2$. 我们可以通过\footnote{这个公式意味着，酉联络 $a=dz\,a_z+d{\bar z}\,a_{\bar z}$，其中 $a_z=-\partial_zf$，具有曲率 $f=da=2\pi i\sum_i n_i\delta_{\rr_i}$. 由具有这一曲率的联络描述的线丛确实是 ${\cal O}(\sum_in_i\rr_i)$.} $a_{\bar z}=\partial_{\bar z}f$ 定义 ${\cal L}'$ 的复结构. 由于 $f$ 中存在奇性，这不是纯规范. 现在可以计算 $\delta a_{\bar z}$，其中 $\delta$ 指的是 ${\cal L}'$ 发生变化时的变分，即各个 $\rr_i$ 的变化. 有
\begin{equation}\label{dunkoff}\delta a_{\bar z}d\bar z=
\sum_id
\rr_i{\partial_{r_i}}a_{\bar z}d\bar z= \sum_i d\rr_i\,d\bar
z\,\partial_{\bar z}\partial_{\rr_i} f=-i\sum_i \pi n_i\,
\delta_{\rr_i}\cdot d\rr_i,\end{equation} 其中导数 $\partial_{\rr_i}a_{\bar z}$ 是利用 $a_{\bar z}$ 在远离各个 $\rr_i$ 处反全纯、且在 $z=\rr_i$ 处具有已知奇性行为这一事实计算的. 与 (\ref{wellwrite}) 中一样，我们将 $d\rr_i$ 视为 $\FF_{SU(N)}$ 上取值于 $TD_{\rr_i}$ 的 1-形式. 在 (\ref{dunkoff}) 中，$d\rr_i$ 与 $D$ 上的 $(1,1)$-形式 $\delta_{\rr_i}$ 缩并，得到 $\FF_{SU(N)}$ 上取值于 $D$ 上 $(0,1)$-形式的 1-形式. 该缩并记为 $\delta_{\rr_i}\cdot d\rr_i$. 现在，当约化到 $\MH$ 时，如同 (\ref{dunko}) 的推导中那样，有 $\delta y \,dz=\sum_a dH_a\,\omega^a$. 利用这一事实和 (\ref{dunkoff})，再回顾 $|d^2z|=idz\wedge d\bar z$，可以计算 (\ref{minto}) 得到
\begin{equation}\label{unkoff} \Omega_I=\sum_{a,i}n_i\, dH_a\wedge
(\omega^a,d\rr_i),\end{equation} 与 (\ref{wellwrite}) 一致.\transnote{这里的归一化是：约定 $|d^2z|=i\,dz\wedge d\bar z$，$\delta_{\rr_i}$ 满足 $\int_D\delta_{\rr_i}=1$；由 $\bar\partial\partial f=i\pi\sum_i n_i\delta_{\rr_i}$ 与 $f\sim\frac12\ln|z-\rr_i|^2$ 可核出 (3.38) 的系数 $-i\pi$（关键是 $\partial_{\rr_i}$ 作用在奇性上等于 $-\partial_z$），再乘 $|d^2z|$ 即得 (3.39) 与 (\ref{wellwrite}) 一致。原文该处脚注说“具有曲率 $f=da=2\pi i\sum_i n_i\delta_{\rr_i}$”，其中 $f$ 其实是标量函数而非曲率：由 $a_z=-\partial_zf$、$a_{\bar z}=\partial_{\bar z}f$ 得 $da=2\bar\partial\partial f=2\pi i\sum_i n_i\delta_{\rr_i}$，故曲率是 $2\pi i\sum_i n_i\delta_{\rr_i}$，与 (3.38) 的系数相差因子 2，两者并不矛盾。（来源：A 级，原文 §3.4.1 式 (3.37)–(3.39) 及其脚注；2026-11-19 核对.）}

